% ****** Start of file aipsamp.tex ******
%
%   This file is part of the AIP files in the AIP distribution for REVTeX 4.
%   Version 4.2a of REVTeX, December 2014
%
%   Copyright (c) 2014 American Institute of Physics.
%
%   See the AIP README file for restrictions and more information.
%
% TeX'ing this file requires that you have AMS-LaTeX 2.0 installed
% as well as the rest of the prerequisites for REVTeX 4.2
%
% It also requires running BibTeX. The commands are as follows:
%
%  1)  latex  aipsamp
%  2)  bibtex aipsamp
%  3)  latex  aipsamp
%  4)  latex  aipsamp
%
% Use this file as a source of example code for your aip document.
% Use the file aiptemplate.tex as a template for your document.
\documentclass[%
 aip,
 jmp,%
 amsmath,amssymb,
%preprint,%
 reprint,%
%author-year,%
%author-numerical,%
]{revtex4-2}

\usepackage{graphicx}% Include figure files
\usepackage{dcolumn}% Align table columns on decimal point
\usepackage{bm}% bold math
%\usepackage[mathlines]{lineno}% Enable numbering of text and display math
%\linenumbers\relax % Commence numbering lines

\begin{document}

\preprint{AIP/123-QED}

\title[Raman Spectroscopy]{Review on Experimental Physics:\\Raman Spectroscopy and Machine Learning - A structural analysis}% Force line breaks with \\
\thanks{Footnote to title of article.}

\author{Khanh Gia Bui}
 %\altaffiliation[Also at ]{Physics Department, XYZ University.}%Lines break automatically or can be forced with \\
\author{Duc Anh Tran}%
\author{Duy Tien Nguyen}%
\affiliation{ 
Vietnam National University, Hanoi University of Science%\\This line break forced with \textbackslash\textbackslash
}%

%\author{C. Author}
% \homepage{http://www.Second.institution.edu/~Charlie.Author.}
%\affiliation{%
%Second institution and/or address%\\This line break forced% with \\
%}%

\date{\today}% It is always \today, today,
             %  but any date may be explicitly specified

\begin{abstract}
Raman spectroscopy provides rich and complex structural information, often via spectrally wide-range, digitally resolved intensity responses of light-matter interactions, through the primary mechanism of \textit{Raman scattering}. With the gradual development of Raman spectroscopy and its expanding application domains, conventional data analysis and model analytics have become largely insufficient. As a result, the application of machine learning and deep learning to the interpretation, processing, and modelling of material properties is an active and significant research topic, with substantial progress and success. However, an open question remains: what \textit{structural information} do ML/AI-assisted or -based Raman spectroscopic systems actually capture during their operation? This paper analyses traditional methods of data processing and model formation from Raman spectral data, alongside the new ML/AI systems developed to improve upon conventional approaches. The bulk of the review is directed toward diagnosing the \textit{structural depth} of current machine learning and artificial intelligence systems --- whether they retain or capture structural information, how such information is inferred, and the interpretability and accessibility of that information. Methods reviewed range from Principal Component Analysis, Spectral Regression, and Support Vector Machines to modern Multilayer Perceptron Networks (MLP), Convolutional Neural Networks (CNN), Attention-based Convolutional Neural Networks, and Graph Neural Networks. The limitations of each model, current approaches, and detailed experimental results are discussed alongside directions for future research.%
\end{abstract}

\keywords{Optical physics, Raman scattering, Raman spectroscopy, Spectroscopy, Machine learning, Deep learning, mathematical modelling}%Use showkeys class option if keyword
                              %display desired
\maketitle

%\begin{quotation}
%The ``lead paragraph'' is encapsulated with the \LaTeX\ 
%\verb+quotation+ environment and is formatted as a single paragraph before the first section heading. 
%(The \verb+quotation+ environment reverts to its usual meaning after the first sectioning command.) 
%Note that numbered references are allowed in the lead paragraph.
%%
%The lead paragraph will only be found in an article being prepared for the journal \textit{Chaos}.
%\end{quotation}

\section{Introduction}\label{sec:intro}

Raman spectroscopy, the spectroscopic method utilizing Raman scattering, discovered by C. V. Raman \cite{raman_new_1928}, is one of the more famous approach in use of modern science and application in experimental or practical fields, often dubbed for its properties as fast, convenient, and \textit{nondestructive inferencing behaviours} that are utilized to probe and enable nondestructive technique, being used in sensitive system or sample, and advancing understanding in multiple fields where such experimental properties are required \cite{Chandra2024RamanBiologicalResearch}. Its utilization can be found in various domains, in inorganic and minerals identification or examination, for arts and archaeology, polymer and emulsions with quantitative inference, identification of different substances for fingerprinting databases of their Raman spectras, in bioengineering \cite{MullerDavidBioEngineering2024} and biomedical \cite{Sharma2025SERSBiomedical,Haugen2026InVivoRaman} application, mostly in monitoring, real-time assessment and control probe. They are also used in drug discovery, drug engineering, forensic, chemical process monitoring, and so on\cite{hollas2004modern,EwenFrontRamanModern}.

\section{Overview of the Theory of Raman Scattering}

There are a few 

\clearpage

\appendix

\section{Appendixes}



\clearpage

%\nocite{*}
\bibliography{aipsamp}% Produces the bibliography via BibTeX.

\end{document}
%
% ****** End of file aipsamp.tex ******
